\section{Conclusion}
As programming becomes increasingly integrated into scientific and educational workflows, understanding how students interact with computational notebooks is of crucial importance. This study evaluated a JupyterLab dashboard extension designed to monitor notebook activity and make the iterative, trial-and-error nature of programming transparent to learners. The evaluation demonstrated that while translating complex, cell-level interaction data into an intuitive user interface presents certain layout and usability challenges, the underlying concept is highly viable. Students perceived the dashboard as very easy to use and reported a low subjective mental workload. This is a key success, as it indicates that the system effectively avoids creating additional cognitive hurdles during programming. Although recognizing the immediate instructional usefulness of these visualizations likely requires more longitudinal engagement, targeted pedagogical guidance, and the integration of concrete problem-solving hints, the dashboard lays an important foundation. By shifting away from purely outcome-based evaluations toward process-oriented analytics, tools like this have the potential to sustainably bridge the gap between the correctness of the final code and the students' individual learning paths.